% =====================================================================
%  Model 1 as implemented: Model 1A / Model 1B, output scale (log10),
%  joint marginal-likelihood fit, and the LF subsample.
%
%  Drop-in for Notes_v1.tex. Suggested insertion point: immediately
%  after \subsection{Correlation modelling} (subsec:correlation-modelling),
%  since it instantiates both rho parametrisations introduced there.
%  Uses only macros and labels already defined in Notes_v1.tex
%  (\nHF, \nLFtenk, \nwl, \litref, eq:design_nested, eq:rho-per-bin,
%  subsec:model1, subsec:correlation-modelling, fig:correlation).
% =====================================================================

\subsection{Model 1 as implemented}
\label{subsec:model1-implementation}

This section records the concrete implementation of Model~1
(Section~\ref{subsec:model1}), which departs from the generic description in
three ways that matter for reproducibility. First, the two parametrisations of
the scaling factor of Section~\ref{subsec:correlation-modelling} are
instantiated as \emph{Model~1A} (one global $\rho$ shared by all bins) and
\emph{Model~1B} (a free per-bin $\rho_j$), and compared by leave-one-out
cross-validation rather than by likelihood-based criteria. Second, each model
is fitted on two output scales --- the raw eclipse depths and their
$\log_{10}$ --- giving four configurations in total. Third, the joint
GP fit does not use the recursive two-step scheme sketched in
Section~\ref{subsec:model1}: the LF process, the discrepancy, and $\rho_j$ are
estimated \emph{jointly} through a single marginal likelihood per bin, with
the $N_L=\nLFtenk$ LF sample entering through a fixed random subsample whose
size is set by an explicit memory and wall-clock budget.

\subsubsection{Output scale: linear versus $\log_{10}$ depths}
\label{subsubsec:model1-logscale}

The simulated eclipse depths are strictly positive (minimum
${\sim}1.3\times10^{-6}$) but span roughly four orders of magnitude across
the training set, with strong right skew: deep molecular bands compress most
samples toward small depths while a few hot configurations dominate the linear
scale. We therefore run every Model~1 configuration twice, once on the raw
depths and once on the elementwise transform
\[
    \tilde y_{i,j} \;=\; \log_{10} y_{i,j},
\]
applied to $Y^{H}$, $Y^{L}$, and the decoupled $\nLFtenk$-point LF set
immediately after loading ($\log_{10}$ rather than $\ln$, following astronomy
convention; positivity makes the transform exact, with no clamping). The
atmospheric inputs $\boldsymbol\theta$ are untouched. Exploratory diagnostics
show that the transform removes the skew and slightly \emph{raises} the HF--LF
correlations of Fig.~\ref{fig:correlation}; consistently, the pooled scaling
estimate below increases from $\hat\rho=0.350$ (linear) to $\hat\rho=0.460$
($\log_{10}$).

Two caveats follow from the nonlinearity of the transform.
(i)~Quantities fitted on $\log_{10}$ depths are \emph{different parameters},
not rescalings of their linear-scale counterparts: the AR(1) relation on the
log scale,
$\tilde f_{H,j}=\rho_j\,\tilde f_{L,j}+\tilde\delta_j$, corresponds in linear
units to the multiplicative power law
\[
    y_{H,j} \;=\; 10^{\tilde\delta_j(\boldsymbol\theta)}\,
                  \bigl(y_{L,j}\bigr)^{\rho_j},
\]
so a log-scale $\rho_j$ acts as an exponent, not a slope, and models on the
two scales may only be compared through back-transformed predictions.
(ii)~For a Gaussian predictive law $\N(\mu,s^2)$ in $\log_{10}$ space, the
back-transform $10^{\mu}$ is the predictive \emph{median} of the implied
log-normal (its mean is $10^{\mu}\exp[(s\ln 10)^2/2]$), while central
predictive intervals map monotonically,
$[\,10^{\mu-1.96s},\,10^{\mu+1.96s}\,]$. All cross-scale comparisons reported
below score $10^{\mu}$ against the held-out HF depths in original units. The
observed spectrum contains negative depths (noise around small signals), so in
the downstream inversion the log-scale surrogate requires masking or a
likelihood formulated in linear units; the simulated training set is
unaffected.

\subsubsection{Joint single-likelihood fit of the per-bin AR(1) model}
\label{subsubsec:model1-jointfit}

For each bin $j$ the two-fidelity model of Section~\ref{subsec:model1} is
implemented by augmenting the input with a fidelity indicator
$t\in\{0,1\}$ ($0=$ LF, $1=$ HF) and endowing the scalar process on the
augmented space with the kernel
\begin{equation}
  k_j\bigl((\boldsymbol\theta,t),(\boldsymbol\theta',t')\bigr)
  \;=\;
  \rho_j^{\,t+t'}\; k_{L,j}(\boldsymbol\theta,\boldsymbol\theta')
  \;+\;
  t\,t'\; k_{\delta,j}(\boldsymbol\theta,\boldsymbol\theta'),
  \label{eq:ar1-augmented-kernel}
\end{equation}
with $k_{L,j}$ and $k_{\delta,j}$ the SE--ARD kernels of
Section~\ref{subsec:model1}. Evaluating
\eqref{eq:ar1-augmented-kernel} on the two fidelity groups reproduces exactly
the $2\times2$ block covariance of the autoregressive prior,
\[
    K_{LL}=k_{L,j},\qquad
    K_{LH}=\rho_j\,k_{L,j},\qquad
    K_{HH}=\rho_j^{2}\,k_{L,j}+k_{\delta,j},
\]
and independent noise levels $\sigma^2_{nL,j}$, $\sigma^2_{nH,j}$ are added on
the diagonal of the LF and HF rows respectively. The LF and HF observations of
bin $j$ are stacked into one training vector and the \emph{entire}
per-bin hyperparameter set
\[
  \Theta_j=\bigl\{\rho_j,\;\sigma^2_{L,j},\;\boldsymbol\ell_{L,j,1:9},\;
    \sigma^2_{\delta,j},\;\boldsymbol\ell_{\delta,j,1:9},\;
    \sigma^2_{nL,j},\;\sigma^2_{nH,j}\bigr\},
  \qquad |\Theta_j| = 23,
\]
is estimated by maximising one joint marginal likelihood
\[
    y_j \sim \N\bigl(\mathbf 0,\; K_j + \Sigma_{\mathrm{noise},j}\bigr),
\]
so that $\rho_j$ is informed simultaneously by the LF field and by the
co-located HF residuals, rather than plugged in from a first-stage fit. This
is the classical co-kriging estimation strategy
(\litref{Kennedy \& O'Hagan (2000), \emph{Biometrika} 87(1), the original
autoregressive multi-fidelity model; Forrester, S\'obester \& Keane (2007),
\emph{Proc.\ R.\ Soc.\ A} 463, for the engineering co-kriging practice}); the
recursive alternative, which fits the LF process first and the discrepancy on
its residuals (\litref{Le Gratiet \& Garnier (2014), \emph{Int.\ J.\
Uncertainty Quantification} 4(5)}), was deliberately not used here, at the
price of the joint covariance solve whose cost is quantified below.

Optimisation is gradient-based (L-BFGS-B) in $\log\Theta_j$, which enforces
positivity of every hyperparameter --- including $\rho_j>0$, consistent with
the uniformly positive empirical slopes of Eq.~\ref{eq:rho-per-bin} --- with
one random restart in addition to the default initialisation. The nine
atmospheric inputs are standardised to zero mean and unit variance using the
moments of the LF design; in these units the box constraints are
$\ell\in[10^{-2},10^{2}]$, $\sigma^2\in[10^{-3},10^{3}]$,
$\rho\in[10^{-3},10^{3}]$, and $\sigma^2_{n}\in[10^{-8},10]$, with
$\ell$ initialised at $1$ and $\sigma^2_{n}$ at $10^{-4}$. The stacked
training vector of each bin is centred and scaled internally before fitting,
and a jitter of $10^{-10}$ is added to the diagonal for numerical stability.
The $m=\nwl$ per-bin fits are mutually independent, so the fitted $\rho_j$,
signal variances, ARD length-scales, log marginal likelihoods, and wall times
are recorded per wavelength.

Prediction of a new HF spectrum evaluates each posterior at
$(\boldsymbol\theta_\ast,\,t=1)$: no LF run at $\boldsymbol\theta_\ast$ is
required, because the augmented kernel \eqref{eq:ar1-augmented-kernel}
propagates the LF information through the cross-covariance
$\rho_j\,k_{L,j}$. The predicted spectrum is the stack of the $m$ per-bin
Gaussian posteriors, with the diagonal cross-bin covariance already discussed
in Section~\ref{subsec:model1}.

\subsubsection{Computational budget and the LF subsample}
\label{subsubsec:model1-subsampling}

\paragraph{Why the exact joint fit is infeasible.}
The exact joint design per bin stacks all LF and HF rows,
$n = N_L + N_H = 10\,097$. Each L-BFGS-B iteration requires (i) a Cholesky
factorisation of the $n\times n$ joint covariance, $\mathcal O(n^3)$ time, and
(ii) the kernel-gradient stack of $n^2\,|\Theta_j|$ doubles needed for the
marginal-likelihood gradient. A direct estimate of the peak memory of one
iteration,
\[
    \bigl(2\,|\Theta_j| + 4\bigr)\, n^2 \times 8 \ \text{bytes}
    \;\approx\; 40.8\ \text{GB}
    \qquad (|\Theta_j|=23,\ n=10\,097),
\]
exceeds the available 17.2\,GB of RAM by more than a factor of two, so the
exact fit cannot even run. Time is equally binding: a measured timing ladder
at $n\in\{497,\,997,\,1\,997\}$ (6.4\,s, 26.6\,s, 219\,s per fit) extrapolated
at the cubic rate gives ${\sim}2.8\times10^{4}$\,s \emph{per bin}, i.e.\
${\sim}1\,536$\,h for all $m=\nwl$ bins. The implementation therefore refuses
the exact fit up front whenever the estimated iteration footprint exceeds
$80\%$ of physical memory, reporting these numbers, and requires any reduction
to be requested explicitly.

\paragraph{Subset-of-data on the LF level only.}
The reduction adopted is the simplest one: a \emph{subset-of-data} (SoD)
approximation in which one uniform random subsample of $n_s$ LF rows is drawn
without replacement (fixed seed, indices recorded with the fitted model) and
stacked with \emph{all} $N_H=\nHF$ HF rows; the same subsample is shared by
all $m$ bins, which keeps the $195$ per-bin fits comparable and lets the joint
design be assembled once. The asymmetry is deliberate. All information about
the discrepancy $\delta_j$ resides in the $\nHF$ co-located pairs
(Eq.~\ref{eq:design_nested}), which are never thinned; the LF field, by
contrast, is abundantly sampled ($\nLFtenk$ points), is the smoother of the
two processes, and --- per the input-sensitivity analysis --- varies along
only ${\sim}3$ of the $9$ input directions, so uniform thinning of the LF
Latin hypercube sacrifices the least information per discarded row while
preserving its marginal space-filling coverage in distribution.

The subsample size is then \emph{maximised against the budget} rather than
chosen by convention: memory alone would admit $n_s\approx5\,800$, but wall
time binds first, and the timing ladder places the largest subsample
compatible with an overnight (${\sim}8$\,h) run at $n_s\approx1\,600$. The
production configuration uses $n_s=1\,700$ LF rows plus the $\nHF$ HF rows per
bin; a reduced setting ($n_s=400$, ${\sim}7$\,s per bin) is retained for fast
iteration during development.

\paragraph{What SoD does and does not guarantee.}
SoD approximates at the \emph{data} level, not the \emph{inference} level:
conditionally on the retained subsample the fitted object is still an exact GP
posterior under the assumed model, so the predictive variances remain valid
--- typically wider, reflecting the discarded points --- rather than
overconfident. This contrasts with structural approximations such as
Nystr\"om or FITC-type inducing-point methods, which modify the prior or
likelihood itself (\litref{Qui\~nonero-Candela \& Rasmussen (2005),
\emph{JMLR} 6, for the unifying view; Rasmussen \& Williams (2006),
\emph{Gaussian Processes for Machine Learning}, Ch.~8.1, for SoD as the
baseline}). Empirically, SoD is a strong per-unit-compute baseline among
sparse GP methods (\litref{Chalupka, Williams \& Murray (2013), \emph{JMLR}
14, a framework for evaluating approximation methods for GP regression}), and
random subsampling admits explicit generalisation-error bounds showing that,
in the data-abundant regime, thinning need not degrade the convergence rate
(\litref{Hayashi, Imaizumi \& Yoshida (2020), AISTATS, on random subsampling
of GP regression}). Variational inducing-point schemes with provably
convergent posteriors (\litref{Titsias (2009), AISTATS; Burt, Rasmussen \& van
der Wilk (2019), ICML, convergence rates for sparse variational GPs}) and
Vecchia-type factorisations (\litref{Vecchia (1988), \emph{JRSS-B} 50;
Katzfuss \& Guinness (2021), \emph{Statist.\ Sci.} 36}) would scale further,
but both replace the exact AR(1) marginal likelihood --- through which
$\rho_j$ is identified jointly --- by a surrogate objective; keeping the
likelihood exact on a reduced design was preferred at the present LF sample
sizes.

\subsubsection{Model 1A versus Model 1B: estimators and cross-validated
comparison}
\label{subsubsec:model1a-vs-1b}

\paragraph{Estimation.}
In the joint GP layer each bin optimises its own $\rho_j$, so the fitted layer
realises Model~1B; constraining a single $\rho$ across bins would couple the
$m$ otherwise-independent likelihoods. The 1A-versus-1B question is therefore
settled upstream, on the closed-form \emph{mean layer} of Model~1 --- the
affine LF$\to$HF map obtained when the discrepancy is held at its prior mean
--- for which both parametrisations have explicit estimators. With per-bin
training moments $\bar y^{L}_{j},\,s_{L,j},\,\bar y^{H}_{j}$, Model~1B uses
exactly the through-origin slope of Eq.~\ref{eq:rho-per-bin},
\begin{equation}
  \hat\rho_j
  \;=\;
  \frac{\operatorname{Cov}_i\!\bigl(y^{H}_{\cdot,j},\,y^{L}_{\cdot,j}\bigr)}
       {\operatorname{Var}_i\!\bigl(y^{L}_{\cdot,j}\bigr)},
\end{equation}
which is the maximum-likelihood estimate of the AR(1) scaling when the
discrepancy is directly observable at the co-located points, while Model~1A
pools the per-bin standardised pairs over all $(i,j)$:
\begin{equation}
  \hat\rho_{\mathrm{A}}
  \;=\;
  \frac{\sum_{j=1}^{m}\sum_{i=1}^{n}
        \tilde r^{H}_{i,j}\,\tilde z^{L}_{i,j}}
       {\sum_{j=1}^{m}\sum_{i=1}^{n}\bigl(\tilde z^{L}_{i,j}\bigr)^{2}},
  \qquad
  \tilde z^{L}_{i,j}=\frac{y^{L}_{i,j}-\bar y^{L}_{j}}{s_{L,j}},
  \quad
  \tilde r^{H}_{i,j}=\frac{y^{H}_{i,j}-\bar y^{H}_{j}}{s_{L,j}} .
  \label{eq:rho-global-pooled}
\end{equation}
Because the per-bin standardisation equalises the denominators
($\sum_i (\tilde z^{L}_{i,j})^2 = n$ for every $j$),
\eqref{eq:rho-global-pooled} coincides with the unweighted mean of the
$\hat\rho_j$; on the full paired data it evaluates to
$\hat\rho_{\mathrm{A}}=0.350$ on the linear scale and $0.460$ on the
$\log_{10}$ scale, while the $\hat\rho_j$ range over $[0.10,\,0.68]$ and
$[0.16,\,0.84]$ respectively --- the wavelength structure of
Fig.~\ref{fig:correlation} that a global scaling cannot express. The
associated held-out predictor,
\begin{equation}
  \hat y_{j}(\boldsymbol\theta_\ast)
  \;=\;
  \bar y^{H}_{j} + \hat\rho_j\,\bigl(y^{L}_{j}(\boldsymbol\theta_\ast)
  - \bar y^{L}_{j}\bigr),
  \label{eq:rho-layer-predictor}
\end{equation}
consumes the paired LF spectrum at the held-out input (available on the nested
design, Eq.~\ref{eq:design_nested}); it is used for model comparison only, not
as the deployed surrogate.

\paragraph{Validation protocol.}
The two parametrisations are compared by leave-one-out cross-validation over
\emph{all} $\nHF$ paired HF samples: each spectrum is predicted by
\eqref{eq:rho-layer-predictor} from a fit that never saw it, with every
training-fold quantity --- the moments
$\bar y^{L}_{j},s_{L,j},\bar y^{H}_{j}$ \emph{and} $\hat\rho$ itself ---
recomputed from the $96$ retained rows, so the held-out predictions are
leakage-free. Both models are scored on identical folds (one fixed seed) by
held-out RMSE, pooled and disaggregated per sample and per wavelength, plus
the paired per-sample RMSE differences. An earlier comparison via AIC and a
nested $F$-test was removed deliberately: effective parameter counts are hard
to define credibly across the closed-form layer, the joint GP layer, and the
forthcoming Model~2 variants, and the $F$ reference distribution assumes
i.i.d.\ residuals that spectrally correlated errors violate, whereas held-out
predictive error is defined identically for every model
(\litref{Arlot \& C\'elisse (2010), \emph{Statist.\ Surv.} 4, survey of
cross-validation for model selection}). The same engine (fit and predict
callbacks around a shared fold assignment) is reused for the GP layers with
$K=5$ folds, where each fold refits the full per-wavelength stack.

\paragraph{Results.}
Table~\ref{tab:model1-loo} reports the four configurations. On both output
scales the per-wavelength scaling wins: Model~1B lowers the pooled held-out
RMSE by ${\sim}3\%$ and wins the paired per-sample comparison on
${\sim}72$--$73\%$ of the $\nHF$ samples, confirming that the structured
per-bin slopes of Eq.~\ref{eq:rho-per-bin} carry genuinely predictive
information rather than noise. The $\log_{10}$ scale helps on top of this:
back-transformed to original units (via the predictive median $10^{\mu}$),
each log-fitted model beats its linear counterpart, and the best
configuration overall is \textbf{Model~1B on $\log_{10}$ depths} (pooled
held-out RMSE $2.72\times10^{-4}$, versus $2.88\times10^{-4}$ for the linear
Model~1A baseline). Model~1B on the $\log_{10}$ scale is therefore the
reference configuration for the joint GP layer and the downstream comparisons.

\begin{table}[t]
  \centering
  \caption{Leave-one-out comparison of the $\rho$ parametrisations
  (Model~1A: global $\rho$; Model~1B: per-bin $\rho_j$) on both output
  scales; all $\nHF$ paired samples held out once, identical folds, seed
  fixed. Native-scale RMSE is in eclipse-depth units (linear) or dex
  ($\log_{10}$); the last column back-transforms the $\log_{10}$ predictions
  (median $10^{\mu}$) to depth units for the cross-scale comparison.
  ``B wins'' is the fraction of samples on which Model~1B's held-out
  per-sample RMSE beats Model~1A's within the same scale.}
  \label{tab:model1-loo}
  \begin{tabular}{llccc}
    \toprule
    Output scale & Model & RMSE (native) & RMSE (depth units) & B wins \\
    \midrule
    linear      & 1A (global $\rho=0.350$) & $2.884\times10^{-4}$ & $2.884\times10^{-4}$ & --- \\
    linear      & 1B ($\rho_j$)            & $2.804\times10^{-4}$ & $2.804\times10^{-4}$ & $73.2\%$ \\
    \midrule
    $\log_{10}$ & 1A (global $\rho=0.460$) & $0.1534$ dex & $2.809\times10^{-4}$ & --- \\
    $\log_{10}$ & 1B ($\rho_j$)            & $0.1483$ dex & $2.723\times10^{-4}$ & $72.2\%$ \\
    \bottomrule
  \end{tabular}
\end{table}
